% =============================================================
% Relativistic Variational Principle for Rayleigh-Taylor (Ch X, §93)
% Agent 43: Relativistic extension of the general variational principle
% =============================================================
%
% Classical source: output/chapters/chapter_10.tex, lines 412--577
% Conventions: RELATIVISTIC_CONVENTIONS.md, rel_preamble.tex
%
% This file extends Chandrasekhar's general variational principle
% for Rayleigh-Taylor instability (§93) to the relativistic regime.
%
% Key physical points:
%   (i)   The inertia density rho is replaced by the relativistic
%         enthalpy density w = (epsilon + p)/c^2 throughout the
%         kinetic-energy functional.
%   (ii)  The gravitational potential energy retains D(rho) but
%         acquires relativistic stratification corrections from the
%         Tolman-Oppenheimer-Volkoff equilibrium.
%   (iii) The viscous dissipation functional generalises to
%         BDNK first-order causal dissipation.
%   (iv)  All results reduce to Chandrasekhar's §93 in the limit
%         c -> infty (where w -> rho).
% =============================================================

\section{Relativistic variational principle for Rayleigh--Taylor
  instability (\S93 -- relativistic)}\label{sec:rel-10-93}

In the classical treatment of \S\,93, Chandrasekhar derived a general
variational principle governing the characteristic values of the
Rayleigh--Taylor problem, including the effects of viscosity and
surface tension.  We now extend this formulation to the relativistic
regime, wherein the fluid inertia is governed by the \emph{enthalpy
density} rather than the rest-mass density alone, and the dissipative
fluxes obey the BDNK first-order causal constitutive relations
\citep{BemficaDN2018,Kovtun2019}.

\subsection{Relativistic preliminaries}

For a relativistic perfect fluid the energy-momentum tensor is
\begin{equation}\label{eq:rel-10-Tmunu}
  T^{\mu\nu} = \frac{\varepsilon + p}{c^{2}}\,u^{\mu}\,u^{\nu}
  + p\,g^{\mu\nu},
\end{equation}
where $\varepsilon$ is the energy density (including rest-mass
energy), $p$ the isotropic pressure, and $u^{\mu}$ the fluid
four-velocity normalised by $u^{\mu}u_{\mu} = -c^{2}$.  The
\emph{relativistic enthalpy density} (or inertia density) is
\begin{equation}\label{eq:rel-10-enthalpy}
  w \;\equiv\; \frac{\varepsilon + p}{c^{2}}.
\end{equation}
In the non-relativistic limit $\varepsilon \to \rho c^{2}$ and
$p/\rho c^{2} \to 0$, so that $w \to \rho$.

The equilibrium of a relativistic stratified layer in a gravitational
field $g$ is governed by the Tolman--Oppenheimer--Volkoff condition
\begin{equation}\label{eq:rel-10-TOV}
  Dp_{0} = -\bigl(\varepsilon_{0} + p_{0}\bigr)\,\frac{g}{c^{2}}
  = -w_{0}\,g,
\end{equation}
where $D = \mathrm{d}/\mathrm{d}z$ and the subscript~0 denotes
equilibrium quantities.

\subsection{Linearised perturbation equations}

Linearising the relativistic Euler equation (with BDNK
viscosity) about this equilibrium, one obtains the perturbation
equations analogous to the classical \eqref{eq:10-70a}
and~\eqref{eq:10-71}:
\begin{align}\label{eq:rel-10-70a}
  D\,\delta\hat{p} &= -n\,w\,\hat{w}
  + \hat{\mu}(D^{2}-k^{2})\hat{w}
  + 2(D\hat{\mu})(D\hat{w})
  + \frac{g}{n}\,(Dw)\,\hat{w}
  \nonumber\\
  &\quad
  - \frac{k^{2}}{n}\sum_{s}\hat{T}_{s}\,\hat{w}_{s}\,
    \delta(z - z_{s})
  - \hat{\mu}\,k^{2}\hat{w}
  + 2(D\hat{\mu})(D\hat{w}),
\end{align}
and
\begin{equation}\label{eq:rel-10-71}
  k^{2}\,\delta\hat{p}
  = -n\,w\,D\hat{w}
  - k^{2}\hat{\mu}\,D\hat{w}
  + k^{2}\hat{w}\,D\hat{\mu}
  + D\!\bigl(\hat{\mu}\,D^{2}\hat{w}\bigr),
\end{equation}
where:
\begin{itemize}
\item $\hat{w}$ is the vertical-velocity perturbation amplitude
  ($\propto e^{nt}$);
\item $\delta\hat{p}$ is the pressure perturbation;
\item $n$ is the growth rate (eigenvalue);
\item $k$ is the horizontal wave number;
\item $w = (\varepsilon + p)/c^{2}$ replaces $\rho$ as the inertia
  density;
\item $Dw$ replaces $D\rho$ in the buoyancy coupling;
\item $\hat{\mu}$ is the relativistic dynamic viscosity; in the
  BDNK framework $\hat{\mu} = \eta$, the shear viscosity, which
  enters the first-order constitutive relation
  $\Pi^{\mu\nu} = -2\eta\,\sigma^{\mu\nu}$ directly (no
  relaxation-time factor);
\item $\hat{T}_{s}$ is the surface tension at interface~$s$.
\end{itemize}
The boundary conditions remain
\begin{equation}\label{eq:rel-10-72}
  \hat{w} = 0 \quad\text{and}\quad D\hat{w} = 0
  \quad\text{on bounding surfaces.}
\end{equation}

\subsection{Orthogonality relations}

Let $n_{i}$, $n_{j}$ be two distinct characteristic values and
$\hat{w}_{i}$, $\hat{w}_{j}$ the corresponding eigenfunctions.
Proceeding exactly as in the classical derivation --- multiplying the
equations for mode~$i$ by $\hat{w}_{j}$, integrating over the layer
$0 \leqslant z \leqslant d$, and integrating by parts --- one arrives
at the relativistic analogue of equation~\eqref{eq:10-81}:
\begin{align}\label{eq:rel-10-81}
  -n_{i}\int_{0}^{d} w\!\left\{\hat{w}_{i}\hat{w}_{j}
  + \frac{1}{k^{2}}(D\hat{w}_{i})(D\hat{w}_{j})\right\}\mathrm{d}z
  &+ \frac{g}{n_{i}}\int_{0}^{d}(Dw)\,\hat{w}_{i}\hat{w}_{j}\,
    \mathrm{d}z
  \nonumber\\
  &- \frac{k^{2}}{n_{i}}\sum_{s}\hat{T}_{s}\,
    \hat{w}_{i}(z_{s})\hat{w}_{j}(z_{s})
  \nonumber\\
  &= \int_{0}^{d}\!\Bigl\{(D^{2}\hat{\mu})\,\hat{w}_{i}\hat{w}_{j}
  + 2\hat{\mu}(D\hat{w}_{i})(D\hat{w}_{j})
  \nonumber\\
  &\quad
  + \frac{1}{k^{2}}\hat{\mu}(D^{2}\hat{w}_{i})
    (D^{2}\hat{w}_{j})\Bigr\}\mathrm{d}z.
\end{align}
The crucial change is $\rho \to w$ in every inertial integral on the
left-hand side.

Interchanging $i$ and $j$ and subtracting, we obtain the orthogonality
relation
\begin{align}\label{eq:rel-10-83}
  (n_{j} - n_{i})\!\Biggl[
  &\int_{0}^{d} w\!\left\{\hat{w}_{i}\hat{w}_{j}
  + \frac{1}{k^{2}}(D\hat{w}_{i})(D\hat{w}_{j})\right\}\mathrm{d}z
  \nonumber\\
  &+ \frac{1}{n_{i}n_{j}}\biggl\{
    g\int_{0}^{d}(Dw)\,\hat{w}_{i}\hat{w}_{j}\,\mathrm{d}z
    - k^{2}\sum_{s}\hat{T}_{s}\,\hat{w}_{i}(z_{s})
      \hat{w}_{j}(z_{s})\biggr\}
  \Biggr] = 0.
\end{align}

\subsection{Reality of characteristic values}

If $n_{i}$ is complex, let $n_{j} = n_{i}^{*}$ (complex conjugate),
$\hat{w}_{j} = \hat{w}_{i}^{*}$.  From~\eqref{eq:rel-10-83}:
\begin{equation}\label{eq:rel-10-86}
  \int_{0}^{d} w\!\left\{|\hat{w}|^{2}
  + \frac{1}{k^{2}}|D\hat{w}|^{2}\right\}\mathrm{d}z
  + \frac{1}{|n|^{2}}\biggl\{
    g\int_{0}^{d}(Dw)\,|\hat{w}|^{2}\,\mathrm{d}z
    - k^{2}\sum_{s}\hat{T}_{s}\,|\hat{w}(z_{s})|^{2}
  \biggr\} = 0.
\end{equation}
Since $w > 0$ everywhere, the first integral is positive definite.  In
the absence of surface tension, if $Dw > 0$ (i.e.\ the heavier fluid
underlies the lighter --- \emph{stable} stratification), the equation
cannot be satisfied, and $n$ must be real.  This is the relativistic
counterpart of Chandrasekhar's classical result.

Similarly, using the analogue of~\eqref{eq:10-85}, one can show that
if $n$ is complex then
\begin{align}\label{eq:rel-10-87}
  2\,\operatorname{Re}(n)\int_{0}^{d}
    w\!\left\{|\hat{w}|^{2}
    + \frac{1}{k^{2}}|D\hat{w}|^{2}\right\}\mathrm{d}z
  &= -\int_{0}^{d}(D^{2}\hat{\mu})\,|\hat{w}|^{2}\,\mathrm{d}z
  \nonumber\\
  &\quad
  - \int_{0}^{d}\hat{\mu}\!\left(k^{2}|\hat{w}|^{2}
    + 2|D\hat{w}|^{2}
    + \frac{1}{k^{2}}|D^{2}\hat{w}|^{2}\right)\mathrm{d}z.
\end{align}
Hence, if $D^{2}\hat{\mu} \geqslant 0$ and $\hat{\mu} > 0$
(guaranteed in the BDNK framework with $\eta > 0$), then
$\operatorname{Re}(n) < 0$ whenever $n$ is complex, so
\emph{overstability cannot occur} --- precisely as in the classical
theory.

\subsection{The relativistic variational principle}\label{ssec:rel-10-var}

Setting $i = j$ in~\eqref{eq:rel-10-81} and dropping subscripts, we
obtain the \emph{relativistic energy balance}:
\begin{align}\label{eq:rel-10-88}
  n\underbrace{\int_{0}^{d}
    w\!\left\{\hat{w}^{2}
    + \frac{1}{k^{2}}(D\hat{w})^{2}\right\}\mathrm{d}z}_{I_{1}}
  &- \frac{g}{n}\underbrace{\int_{0}^{d}
    (Dw)\,\hat{w}^{2}\,\mathrm{d}z}_{I_{2}}
  + \frac{k^{2}}{n}\sum_{s}\hat{T}_{s}\hat{w}_{s}^{2}
  \nonumber\\
  &= -\underbrace{\int_{0}^{d}\!\left\{
    \hat{\mu}\!\left[k^{2}\hat{w}^{2}
    + 2(D\hat{w})^{2}
    + \frac{1}{k^{2}}(D^{2}\hat{w})^{2}\right]
    + (D^{2}\hat{\mu})\,\hat{w}^{2}
  \right\}\mathrm{d}z}_{I_{3}}.
\end{align}

\medskip\noindent\textbf{Physical interpretation.}
\begin{itemize}
\item $I_{1}$ is the \textbf{relativistic kinetic energy} functional.
  The replacement $\rho \to w = (\varepsilon + p)/c^{2}$ reflects the
  fact that in relativity both internal energy and pressure contribute
  to the inertia.  Explicitly:
  \begin{equation}\label{eq:rel-10-KE}
    \mathcal{T} \;=\; \tfrac{1}{2}\,
    \frac{\varepsilon + p}{c^{2}}\,|\boldsymbol{v}|^{2}
    \;=\; \tfrac{1}{2}\,w\,|\boldsymbol{v}|^{2},
  \end{equation}
  which reduces to $\tfrac{1}{2}\rho|\boldsymbol{v}|^{2}$ in the
  Newtonian limit.

\item $I_{2}$ is the \textbf{potential energy} from gravitational
  stratification, governed by $Dw$ (the gradient of enthalpy density)
  rather than $D\rho$.  In the Newtonian limit, the
  TOV~\eqref{eq:rel-10-TOV} reduces to $Dp_{0} = -\rho g$, and $Dw
  \to D\rho$, recovering the classical potential-energy functional.

\item $I_{3}$ is the \textbf{viscous dissipation} functional.  In the
  BDNK framework, the viscosity $\hat{\mu} = \eta$ is the shear
  viscosity entering the first-order constitutive relation
  $\Pi^{\mu\nu} = -2\eta\,\sigma^{\mu\nu}$; it is independent of
  the growth rate $n$, so that the eigenvalue problem remains a
  simple quadratic in~$n$.
\end{itemize}

\medskip\noindent\textbf{Variational formulation.}
Consider the effect of an arbitrary variation $\delta\hat{w}$
compatible with the boundary conditions~\eqref{eq:rel-10-72}.  To
first order,
\begin{equation}\label{eq:rel-10-89}
  -\left(I_{1} + \frac{1}{n^{2}}\Bigl[g\,I_{2}
  - k^{2}\sum_{s}\hat{T}_{s}\hat{w}_{s}^{2}\Bigr]\right)\delta n
  = n\,\delta I_{1}
  - \frac{1}{n}\Bigl[g\,\delta I_{2}
    - 2k^{2}\sum_{s}\hat{T}_{s}\hat{w}_{s}\,\delta\hat{w}_{s}\Bigr]
  + \delta I_{3},
\end{equation}
where $\delta I_{1}$, $\delta I_{2}$, $\delta I_{3}$ are the
first-order variations:
\begin{align}
  \tfrac{1}{2}\,\delta I_{1}
  &= \int_{0}^{d}\delta\hat{w}\!\left\{
    w\,\hat{w}
    - \frac{1}{k^{2}}\,D(w\,D\hat{w})\right\}\mathrm{d}z,
  \label{eq:rel-10-91}\\[6pt]
  \tfrac{1}{2}\,\delta I_{2}
  &= \int_{0}^{d}\delta\hat{w}\,(Dw)\,\hat{w}\;\mathrm{d}z,
  \label{eq:rel-10-92}\\[6pt]
  \tfrac{1}{2}\,\delta I_{3}
  &= \int_{0}^{d}\delta\hat{w}\!\left\{
    k^{2}\hat{\mu}\,\hat{w}
    - 2D(\hat{\mu}\,D\hat{w})
    + (D^{2}\hat{\mu})\,\hat{w}
    + \frac{1}{k^{2}}D^{2}(\hat{\mu}\,D^{2}\hat{w})\right\}
    \mathrm{d}z.
  \label{eq:rel-10-93}
\end{align}
Combining these in accordance with~\eqref{eq:rel-10-89}, the
condition $\delta n = 0$ for all admissible $\delta\hat{w}$ becomes
\begin{align}\label{eq:rel-10-94}
  \tfrac{1}{2}\,k^{2}\Bigl\{I_{1}
  &+ \frac{1}{n^{2}}\bigl[g\,I_{2}
    - k^{2}\sum_{s}\hat{T}_{s}\hat{w}_{s}^{2}\bigr]\Bigr\}
    \,\delta n
  \nonumber\\
  &= \int_{0}^{d}\delta\hat{w}\,\biggl\{
    D\!\left[w\,D\hat{w}
      - \frac{\hat{\mu}}{n}(D^{2} - k^{2})D\hat{w}
      - \frac{1}{n}(D\hat{\mu})(D^{2}+k^{2})\hat{w}
    \right]
  \nonumber\\
  &\qquad
    - k^{2}\!\left[
      w\,\hat{w}
      - \frac{\hat{\mu}}{n}(D^{2} - k^{2})\hat{w}
      - \frac{g}{n^{2}}(Dw)\,\hat{w}
      + \frac{k^{2}}{n^{2}}\sum_{s}\hat{T}_{s}\,
        \delta(z - z_{s})\hat{w}
    \right]\biggr\}\mathrm{d}z.
\end{align}
The integrand on the right-hand side vanishes identically when
$\hat{w}$ satisfies the relativistic eigenvalue problem.  Therefore,
$\delta n = 0$ for all admissible variations $\delta\hat{w}$ if and
only if $\hat{w}$ is a solution of the characteristic-value problem.
This establishes the \emph{relativistic variational principle}.

\subsection{Eigenvalue problem for the growth rate}

The extremal condition defines the eigenvalue problem: find $n$ and
$\hat{w}$ (not identically zero, satisfying the boundary conditions)
such that~\eqref{eq:rel-10-88} is satisfied.  Rearranging:
\begin{equation}\label{eq:rel-10-eigenvalue}
  \boxed{
  n^{2}\,I_{1}[w;\hat{w}]
  + n\,I_{3}[\hat{\mu};\hat{w}]
  - g\,I_{2}[Dw;\hat{w}]
  + k^{2}\sum_{s}\hat{T}_{s}\hat{w}_{s}^{2}
  = 0,
  }
\end{equation}
which is a quadratic in $n$ (for given $\hat{w}$) since in the BDNK
framework $\hat{\mu} = \eta$ is independent of $n$.  Solving:
\begin{equation}\label{eq:rel-10-nquad}
  n = \frac{-I_{3} \pm \sqrt{I_{3}^{2}
    + 4\,I_{1}\bigl(g\,I_{2}
    - k^{2}\sum_{s}\hat{T}_{s}\hat{w}_{s}^{2}\bigr)}}
    {2\,I_{1}}.
\end{equation}
\emph{Instability} ($n > 0$) requires
\begin{equation}\label{eq:rel-10-instab-cond}
  g\,I_{2} - k^{2}\sum_{s}\hat{T}_{s}\hat{w}_{s}^{2} > 0
  \qquad\Longleftrightarrow\qquad
  g\!\int_{0}^{d}(Dw)\,\hat{w}^{2}\,\mathrm{d}z
  > k^{2}\sum_{s}\hat{T}_{s}\hat{w}_{s}^{2},
\end{equation}
i.e.\ the destabilising effect of adverse enthalpy stratification
($Dw < 0$ for heavy-over-light) must overcome the stabilising effect
of surface tension.

\subsection{Bounds on the growth rate}

From the variational formulation one can extract useful bounds.

\subsubsection*{Upper bound (inviscid case)}

In the inviscid limit ($\hat{\mu} \to 0$, hence $I_{3} \to 0$),
equation~\eqref{eq:rel-10-nquad} gives
\begin{equation}\label{eq:rel-10-ninviscid}
  n^{2} = \frac{g\,I_{2} - k^{2}\sum_{s}\hat{T}_{s}\hat{w}_{s}^{2}}
    {I_{1}}.
\end{equation}
Since the variational principle yields the \emph{maximum} of $n^{2}$
over all admissible trial functions, any trial function provides an
upper bound:
\begin{equation}\label{eq:rel-10-upper-bound}
  n_{\max}^{2} \;\leqslant\;
  \sup_{\hat{w}}\;
  \frac{g\displaystyle\int_{0}^{d}(Dw)\,\hat{w}^{2}\,\mathrm{d}z
    - k^{2}\sum_{s}\hat{T}_{s}\hat{w}_{s}^{2}}
    {\displaystyle\int_{0}^{d}
    w\!\left\{\hat{w}^{2}
    + \frac{1}{k^{2}}(D\hat{w})^{2}\right\}\mathrm{d}z}.
\end{equation}
Compared with the classical bound (in which $w \to \rho$ and $Dw \to
D\rho$), the relativistic bound differs in two respects:
\begin{enumerate}
\item The inertia is \emph{larger}: $w = (\varepsilon + p)/c^{2}
  \geqslant \rho$, so the denominator is larger and the growth rate is
  \emph{smaller}.
\item The driving gradient involves $Dw$ rather than $D\rho$.  Since
  the relativistic equilibrium condition~\eqref{eq:rel-10-TOV}
  implies $Dw = D\rho + D(p/c^{2})$, the effective Atwood number is
  modified.
\end{enumerate}

\subsubsection*{Lower bound (viscous damping)}

For $\hat{\mu} > 0$, the viscous term $I_{3}$ is positive definite
and acts to reduce $n$.  From~\eqref{eq:rel-10-nquad}:
\begin{equation}\label{eq:rel-10-viscous-bound}
  n \;\leqslant\;
  \frac{g\,I_{2} - k^{2}\sum_{s}\hat{T}_{s}\hat{w}_{s}^{2}}
    {I_{3}},
\end{equation}
which places a lower bound on the stabilising effect of viscosity.
In the BDNK framework, the viscosity $\hat{\mu} = \eta$ is the
physical shear viscosity constrained by frame-dependent causal
bounds \citep{BemficaDN2018,Kovtun2019}; it does not depend on the
growth rate~$n$.  The viscous bound~\eqref{eq:rel-10-viscous-bound}
therefore applies directly, without the frequency-dependent
reduction present in second-order (Israel--Stewart) theories.

\subsection{Non-relativistic limit}\label{ssec:rel-10-NR}

The non-relativistic limit is obtained by taking $c \to \infty$
(equivalently, $p/\varepsilon \to 0$; BDNK frame corrections
vanish at $\mathcal{O}(v^2/c^2)$):
\begin{enumerate}
\item $w = (\varepsilon + p)/c^{2} \to \rho c^{2}/c^{2} = \rho$.
\item $Dw \to D\rho$.
\item The TOV condition~\eqref{eq:rel-10-TOV} reduces to $Dp_{0} =
  -\rho g$.
\item The viscosity $\hat{\mu} = \eta
  \to \mu$ (standard Navier--Stokes viscosity).
\item $I_{1} \to \int_{0}^{d}\rho\bigl\{\hat{w}^{2} +
  k^{-2}(D\hat{w})^{2}\bigr\}\,\mathrm{d}z$, $I_{2} \to
  \int_{0}^{d}(D\rho)\,\hat{w}^{2}\,\mathrm{d}z$.
\end{enumerate}
Under these replacements, equations~\eqref{eq:rel-10-88}
through~\eqref{eq:rel-10-94} reduce identically to
Chandrasekhar's equations~\eqref{eq:10-88}--\eqref{eq:10-94}.
The variational principle, eigenvalue equation, and all bounds
recover their classical forms exactly.

\subsection{Summary}

The relativistic variational principle for Rayleigh--Taylor
instability is structurally identical to the classical one, with the
systematic replacement $\rho \to w = (\varepsilon + p)/c^{2}$ in all
inertial quantities and $D\rho \to Dw$ in the buoyancy term.  The
key new features are:
\begin{enumerate}
\item \textbf{Enhanced inertia:} pressure contributes to the effective
  mass density, generally reducing growth rates.
\item \textbf{Modified Atwood number:} the density contrast is
  replaced by the enthalpy-density contrast.
\item \textbf{Causal dissipation:} the BDNK first-order framework
  provides causal viscous damping with viscosity $\eta$ subject to
  frame-dependent bounds; no relaxation-time weakening occurs.
\item \textbf{Causality bound:} growth rates are bounded by $c/d$,
  since perturbations cannot propagate superluminally; this is
  guaranteed by the BDNK strong hyperbolicity
  \citep{BemficaDN2018,Kovtun2019}.
\end{enumerate}

%% ================================================================
\subsection{Astrophysical application: Variational bounds on RT growth
rates}\label{ssec:rel-10-93-astro}

The variational principle provides a practical tool for estimating
maximum RT growth rates in astrophysical contexts where exact
solutions are unavailable.  For a relativistic exponentially stratified
fluid layer (as may arise in neutron star envelopes or post-merger
remnants), the upper bound on the growth rate~$n$ depends on both
the wavenumber and the relativistic parameter
$\xi = p/(\rho c^{2})$ characterising the equation of state
\citep{RezzollaZanotti2013}.

The key result is that the relativistic inertia enhancement
$w = \rho c^{2}(1+\xi)$ systematically \emph{lowers} the variational
upper bound on $n^{2}$ compared with the classical prediction, since
the denominator $I_{1}[w;\hat{w}]$ in the Rayleigh
quotient~\eqref{eq:rel-10-upper-bound} increases with~$\xi$.  For
a radiation-dominated post-shock layer ($\xi = 1/3$), the maximum
growth rate is reduced by a factor $\sim 1/\sqrt{1+\xi} \approx 0.87$
relative to the Newtonian estimate \citep{Radice2016}.

Figure~\ref{fig:variational-bound-RT} illustrates the variational
bounds for the exponentially stratified case.

\begin{figure}[htbp]
\centering
\includegraphics[width=0.95\textwidth]{fig_variational_bound_RT.pdf}
\caption{Variational bounds on the RT growth rate for an exponentially
  stratified relativistic fluid.  \textit{Left:} Exact growth rates
  for modes $m = 1, 2, 3$ (solid) compared with the upper bound from
  a parabolic trial function (dotted).  The asymptotic limit
  $n_{\max} = \sqrt{g\beta}$ is shown as a dashed line.
  \textit{Right:} Maximum growth rate ($m=1$) for different values
  of the relativistic parameter~$\xi$, showing the systematic
  reduction of the bound due to enhanced inertia.}
\label{fig:variational-bound-RT}
\end{figure}
